% try: change a bond in the caffeine molecule, or balance another reaction with \ce{...}
\documentclass[11pt]{article}
\usepackage[margin=1in]{geometry}
\usepackage{chemfig}
\usepackage{amsmath}
\usepackage[version=4]{mhchem}
\setchemfig{atom sep=2em, bond style={line width=.8pt}}

\begin{document}
\begin{center}{\LARGE\bfseries Chemistry}\\[3pt]{\large molecules and reactions, typeset}\end{center}

\section*{Molecules}
\begin{center}
\begin{tabular}{ccc}
  \chemfig{*6((=O)-N(-CH_3)-*5(-N=-N(-CH_3)-=)--(=O)-N(-H_3C)-)} &
  \chemfig{*6(-=-=-=)} &
  \chemfig{H_3C-CH_2-OH} \\[6pt]
  caffeine & benzene & ethanol
\end{tabular}
\end{center}

\section*{A reaction}
\begin{center}
\schemestart
  \chemfig{*6(-=-=-=)}
  \arrow{->[\ce{HNO3}][\ce{H2SO4}]}
  \chemfig{*6(-=-(-NO_2)=-=)}
\schemestop
\end{center}

\section*{Equations}
\begin{align*}
  &\ce{2 H2 + O2 -> 2 H2O} \\
  &\ce{CO2 + H2O <=> H2CO3 <=> H+ + HCO3-} \\
  &\ce{CaCO3 ->[\Delta] CaO + CO2 ^} \\
  &\ce{Cu^{2+} + 4 NH3 -> [Cu(NH3)4]^{2+}}
\end{align*}

\end{document}
